% Generated draft: requires mathematical and source review.
\documentclass[11pt]{article}
\usepackage[margin=25mm]{geometry}
\usepackage{fontspec}
\setmainfont{Latin Modern Roman}
\usepackage{amsmath,amssymb,mathtools,unicode-math}
\setmathfont{Latin Modern Math}
\usepackage{enumitem,needspace,xurl,hyperref}
\setlength{\parindent}{0pt}
\setlength{\parskip}{5pt}
\setlength{\emergencystretch}{3em}
\newcommand{\N}{\mathbb{N}}
\newcommand{\Z}{\mathbb{Z}}
\newcommand{\Q}{\mathbb{Q}}
\newcommand{\R}{\mathbb{R}}
\newcommand{\C}{\mathbb{C}}
\newcommand{\F}{\mathbb{F}}
\newcommand{\A}{\mathbb{A}}
\newcommand{\PP}{\mathbb{P}}
\newcommand{\E}{\mathbb{E}}
\newcommand{\eps}{\varepsilon}
\newcommand{\dd}{\,\mathrm{d}}
\begin{document}
\section*{UnsolvedMath ID: 30003277}
\textbf{Source problem code:} catalogue entry
\subsection{Definitions and mathematical statement}
\begin{enumerate}
\item \textbf{A Dirichlet series with one zero.} For a real number \(\alpha\), write
\[
H_\alpha=\{s\in\C:\operatorname{Re}s>\alpha\}.
\]
In the ordinary Dirichlet-series interpretation, a Dirichlet series has the form
\[
D(s)=\sum_{n=1}^{\infty}a_n n^{-s},\qquad a_n\in\C.
\]
The question is whether there exist such coefficients and \(\alpha\in\R\) for which the series converges throughout \(H_\alpha\) and
\[
\#\{s\in H_\alpha:D(s)=0\}=1.
\]
This displayed condition counts distinct zero locations; the source does not specify whether the unique zero must also be simple. The ordinary-series interpretation and unrestricted complex coefficients are conventions for displaying the question, not hypotheses confirmed by the supplied record. The record does not clarify whether a more restricted class of Dirichlet series is intended. Its original wording additionally says ``unconditional Dirichlet series,'' whose meaning is not explained; in particular, it is unclear whether this concerns unconditional convergence or a construction requiring no unproved hypotheses. These qualifications remain unresolved.

\item \textbf{Gaps between Maass eigenvalues.} Let
\[
\mathbb H=\{x+iy\in\C:y>0\},\qquad
X=\operatorname{SL}_2(\Z)\backslash\mathbb H,
\]
where the group acts by fractional linear transformations. Equip \(\mathbb H\) with its hyperbolic metric; the customary nonnegative Laplace operator is
\[
\Delta=-y^2\left(\frac{\partial^2}{\partial x^2}+\frac{\partial^2}{\partial y^2}\right).
\]
For the intended sequence of positive Laplace eigenvalues of Maass forms on \(X\), listed as
\[
0<\lambda_1\leq\lambda_2\leq\cdots,
\]
the source asks whether both assertions hold:
\[
\liminf_{j\to\infty}(\lambda_{j+1}-\lambda_j)=0,
\qquad
\limsup_{j\to\infty}(\lambda_{j+1}-\lambda_j)=+\infty.
\]
The record does not specify the precise class of Maass forms, whether eigenvalues are repeated according to multiplicity, or whether parity or other restrictions are imposed. Thus the spectral sequence must be clarified before these become fully specified propositions. The displayed ordering concerns a discrete sequence and does not by itself specify how the continuous spectrum is excluded.

\item \textbf{Quadratic roots in short intervals.} For an integer modulus \(m\geq1\), an integer \(a\), and a bounded real interval \(I\), define
\[
N(m,a;I)=\#\{x\in\Z\cap I:x^2\equiv a\pmod m\},
\]
and define the total number of root classes by
\[
R(m,a)=\#\{\bar x\in\Z/m\Z:\bar x^{\,2}=\bar a\}.
\]
The endpoint choices are part of the specified interval \(I\). The source asks whether one can uniformly bound \(N(m,a;I)\) for intervals of length
\[
|I|=m^{1-\varepsilon},
\]
when \(m\) is highly composite and the congruence has many roots. In its standard technical meaning, ``highly composite'' means
\[
\tau(m)>\tau(n)\quad\text{for every integer }1\leq n<m,
\qquad
\tau(n)=\#\{d\geq1:d\mid n\}.
\]
The record does not establish that this technical meaning is intended. It also gives no quantified meaning of ``many,'' no range or quantification for \(\varepsilon\), no assumptions on \(a\) such as coprimality with \(m\), and no proposed bound or permitted dependence of its constants. Consequently, this part specifies a counting problem but does not determine a unique uniform-bound conjecture. The interval convention above makes the counting function explicit; the source itself does not specify endpoint conventions.
\end{enumerate}

The three questions above follow the supplied corrected statement. The supplied verification note says that two clauses appearing in the original statement were removed as duplicates of \texttt{OWR-15174-011}; they are not additional targets of this corrected record.
\subsection{Short English statement}
Permanent UnsolvedMath ID \(30003277\) records three questions: whether a Dirichlet series can have exactly one zero in a half-plane of convergence, whether consecutive Maass eigenvalue gaps have lower limit zero and unbounded upper limit, and how quadratic congruence roots concentrate in short intervals. Several hypotheses and quantifiers remain unspecified. The supplied open status and literature assessment are dated curation, not independently verified conclusions.
\subsection{Sources}
\begin{itemize}
\item Local UnsolvedMath registry: supplied source record for permanent ID \(30003277\), source problem code \texttt{catalogue entry}, titled ``Three Problems from an Analytic Number Theory Session,'' in the Oberwolfach Reports Open Problems collection. Repository registry locations are \texttt{data/unsolvedmath/problems.json} and \texttt{catalog/problems.json}; this draft uses the supplied record without independently inspecting either file. No individual proposer is supplied; the proposed year is \(2016\).
\item Supplied report citation: \emph{Analytic Number Theory (2017)}, Oberwolfach Reports, pp.\ 2975--3030, DOI \texttt{10.4171/owr/2016/53}. Source URL and literature-source link labelled ``OWR 53 report'': \url{https://doi.org/10.4171/owr/2016/53}.
\item Supplied literature-source link labelled ``Withdrawn Hilberdink--Saias preprint'': \url{https://arxiv.org/abs/1812.11880}. The supplied assessment says that the proposed Dirichlet-series construction is not established and that this later preprint was withdrawn because of an unrepairable error in its main proof. This assessment is reproduced as supplied attribution, without independently verifying the preprint or withdrawal claim.
\item Literature checked at: \texttt{2026-08-21}. This exact date identifies supplied dated curation, not an independently verified literature assessment. The background labels the status ``open'' and explicitly requires independent verification of time-sensitive status and cited claims; its assessment specifically discusses the Dirichlet-series question.
\item The supplied record labels its statement \texttt{corrected\_verified}. Its verification note reports that the official problem-session report was checked and that two clauses duplicated in \texttt{OWR-15174-011} were removed, leaving three distinct questions. These are supplied verification claims, not checks performed for this draft.
\end{itemize}
\end{document}
